\section[Specific constructions]{Specific constructions (Besondere Aufbauten)} \label{sec:despco}
\subsubsection{Blackletter script (Fraktur)}
The German blackletter script, or \bt{Fraktur} is an old $16^{th}$ century script that is often found in equally old texts,  on street signs in \bt{Deutschland}, or in traditional businesses. \bt{Fraktur} is no longer used and has been phased out, but has its own angular charm. If one desires, one may seek:
\begin{example}
    \begin{tabular}{m{7cm}|m{7cm}}
    $\mathfrak{Das}$ $\mathfrak{ist}$ $\mathfrak{Fraktur}.$ $\mathfrak{Aber}$ $\mathfrak{ob}$ $\mathfrak{die}$ $\mathfrak{Schrift}$ $\mathfrak{noch}$ $\mathfrak{es}$ $\mathfrak{wert}$ $\mathfrak{ist?}$ & Das ist Fraktur. Aber ob die Schrift noch es wert ist?
    \end{tabular}
\end{example}
\subsubsection{Case Mirroring (Fallspiegelung)}
When \bt{als, wie} are used as comparators, the case of their second argument mirrors that of their first:
\begin{example}
    \begin{tabular}{m{7cm}|m{7cm}}
        \bt{Er($1$) war so groß wie sie($1$)} & He was as big as her\\\hline
        \bt{Besser dem Kind($4$) etwas schenken als dem Mann($4$)} & Better to the child \ut{something gift} than to the man
    \end{tabular}
\end{example}
English does  this for all cases except $1$.
\subsubsection{Odd words (Seltsame Wörter)}
Some adjectives  look like \bt{MdV} or \bt{MdG} forms of  verbs that don't exist:
\begin{example}
    \begin{tabular}{c c c|c c c}
        \multicolumn{3}{c}{\textit{MdV}} & \multicolumn{3}{|c}{\textit{MdG}}\\\hline
        verzwickt & berühmt & berüchtigt & anwesend & weitgehend & abwesend
    \end{tabular}
\end{example}
The meanings of some \bt{MdV} and \bt{MdG} adjectives have  no obvious relation  to the verbs they are derived from:
\begin{example}
    \begin{tabular}{c c|c c}
        \multicolumn{2}{c|}{\textit{MdV}} & \multicolumn{2}{c}{\textit{MdG}}\\\hline
        gelassen & \bt{serene} & dringend & \bt{urgent}\\\hline
        besessen & \bt{obsessed} & umgehend & \bt{immediately}
    \end{tabular}
\end{example}
The superlative phrase `One of the most\ldots' can take any general adjective, or even stay without it:
\begin{example}
    \begin{tabular}{m{14cm}}
    Einer der großen/größeren/größesten Schriftsteller Deutschlands\ldots\\\hline
    Eines seiner Versehen ist dafür Schuld
    \end{tabular}
\end{example}
\subsection[Forced/Permitted actions]{Gezwungene/Gestattete Handlungen}
Marker verbs are employed for both types:
\begin{setdef}
        $\vm{D} = \begin{bmatrix} \texttt{lassen} & \texttt{lassen} \end{bmatrix}$
\end{setdef}
The distinction between forced and permitted actions occurs via context or reformulation, since both are marked by \bt{lassen}. The diligent reader will question immediately the fact that no $E\in\vm{D}$ can $\in\mathcal{S}$, which contradicts the definition made in \autoref{sec:deverb}. This objection, while correct, does not affect any of the explanations made until this point.

\bt{Lassen} has only three possible use cases:
\begin{enumerate}
    \item Main verb without a verb argument
    \item Passive replacement: \bt{sich lassen}
    \item Forced/Permitted actions, which need a verb argument
\end{enumerate}
The irregularities that deny any $E$ from having complete membership of $\mathcal{S}$ lie solely in the passive voice. As a main verb without an argument, everything works as normal, while passive replacements cannot be subject to another passive voice change. Hence, the only exceptional configurations exhibited by \bt{lassen} occur for forced/permitted actions. Add to this the fact that such sentences normally do not employ the passive voice,  makes $\texttt{lassen}\in\mathcal{S}$ hold for most practical applications.

The sole irregularity is that \bt{lassen} uses its \bt{MdV} form when passived, but not when \bt{Perfekt}ed if it takes a verb argument:
\begin{example}
    \begin{tabular}{m{6.5cm}|m{6.5cm}}
        \bt{Ich muss mir meine Weisheitszähne ziehen lassen} & I must \ut{from myself my wisdom teeth pull let}\\\hline
        \bt{Jemand hat seine Münze fallen lassen} & Someone has his coin \ut{fall let}\\\hline
        \bt{Ich werde lernen gelassen} & I am being made to learn\\\hline
        \bt{Er lässt mich lernen} & He is making me learn\\\hline
        \bt{Lass mich gehen!} & Let me go!
    \end{tabular}
\end{example}
\subsection[Cases Akkusativ, Dativ]{Fälle Akkusativ, Dativ}
There exists a unique interaction between these cases when motion of any sort is involved, which has been hinted at in \autoref{sec:deprep} while explaining the forms: \bt{hinter/hinten, unter/unten\ldots}. The destination of any such movement  always has an associated preposition, and stays in \bt{Akk}($2$), assuming that the  preposition supports this. If not, the required case of the preposition takes over.

If a location of rest is declared without motion, \bt{Dat}($7$) is used for it. Any time phrases always use \bt{Dat}($7$) irrespective of motion:
\begin{example}
    \begin{tabular}{m{7cm}|m{7cm}}
        \textit{Akkusativ} & \textit{Dativ}\\\hline
        Sie hängt den Mantel an die Wand & Der Mantel hängt an der Wand\\\hline
        Sie fährt in die Stadt & Sie fährt nach der Stadt\\\hline
        \multicolumn{2}{c}{Zeit Dativ}\\\hline
        Regen beginnt in zwei Stunden & Ich bin vor vier Jahren angekommen
    \end{tabular}
\end{example}
Although this rule has no exceptions,  it may sometimes seem like $7$  is used despite possible interpreted motion:
\begin{example}
    Wir treffen uns am(an + dem) Bahnhof
\end{example}
This is not an exception, since the movement target is assumed as the other person, not the \bt{Bahnhof} itself.
\subsection[Contractions]{Verkürzungen}
These contractions are different from the ones considered earlier, since there is no reformulation  involved, and are merely methods of compressing  words  for ease of speech:
\begin{itemize}
\item Some articles can be pushed into prepositions
\item If the contracted word is \bt{es}, the \textit{EoL}(') is  employed
\item The \bt{ei} at the beginning of any indefinite article can be replaced by the \textit{EoL}
\end{itemize}
\begin{setdef}
    \begin{tabular}{c c|c c}
        für/das & fürs & zu/dem & zum*\\\hline
        zu/der & zur & an/dem & am\\\hline
        an/das & ans & auf/das & aufs\\\hline
        von/dem & vom & hinter/dem & hinterm\\\hline\hline
        \multicolumn{2}{c}{Los geht es!} & \multicolumn{2}{|c}{Los geht's!}\\\hline\hline
        einen & 'nen & einem & 'nem
    \end{tabular}
\end{setdef}
\bt{Zum*} and \bt{zudem} mean entirely different things, which makes this a case of forced contraction. Any other contraction is otherwise not absolutely necessary.
Not every article and preposition can be pushed together, which is why  the above table has been made comprehensive with respect to this rule.

\ut{Fun fact}: Some prepositions like \bt{außer}(apart from)($5$) were used so many times with neutral/male noun singulars, that \bt{außerdem}(in addition to) became a conjunction over time.
\subsubsection{Word elimination (Wortbeseitigung)}
In two specific cases, it is possible to eliminate a word and form a  unit from what remains:
\begin{enumerate}
    \item \bt{Sein} (as the CV for $\vm{T}(1,2)$ in the \bt{Zustandsleidform}) : \bt{Betreten (ist) verboten}
    \item \bt{Von : Eine Menge (von) Leute}
\end{enumerate}
The unit formed by method 1 cannot be used in sentences, and must stand alone. In contrast, method 2 forms a typical noun unit.
\subsubsection{Allgemeine Nennwortverkürzungen mittels Hakenwörter}
Hook words ($\vm{Q}$) can  replace nouns if their previous context is known, and emphasize their connection with any consequent usage. Such an emphasis attempts  to clarify the link between noun units in a paragraph, especially useful if multiple similar nouns are dealt with at once. An example for this type of replacement has been provided in the subsection `Describer subclause' with technical German.
\begin{example}
    \begin{tabular}{m{14cm}}
        Der Mann mag klein sein, aber [dessen] Wirkung ist groß\\\hline
        Mein Vater ist verschwunden. Ich kann [den] nicht mehr sehen\\\hline
        Der Gefangene ist zäh. [Dem] entkommt keine Auskunft
    \end{tabular}
\end{example}
Every such use for $\vm{Q}(1\leftrightarrow3,:)$ has an equivalent  $\vm{P}[a,In,3_p](1\leftrightarrow3,:)$ that can replace it. For $\vm{Q}_i(4,j)$, the equivalence is with $\vm{P}_i[a,De,3_p\{j\}](:,:)$, for any one choice of $i\in\{M,F,N\}$ and $j\in\{1,\mathbb{M}\}$ at a time.
\subsection[Infinitive units]{Grundformeinheiten}
$\vm{T}(:,4)$ stays at the end, with no replacement forms necessary  for any $V\in\mathcal{S}$:
\begin{example}
    \begin{tabular}{m{7cm}|m{7cm}}
        \bt{[Die vermissten Kinder zu suchen] hat  Vorrang} & \bt{*}[To be searching for the missing children] \ut{has} priority\\\hline
        \bt{[Ohne Wasser Arbeit gesucht zu haben] ist  ungewöhnlich} & [To have searched for work without water] is strange\\\hline
        \bt{[In dieser Welt leben zu müssen] wirkt manchmal nervig} & [To have to live in this world] is sometimes irritating\\\hline
        \bt{[Nicht mehr gepflegt zu sein] kann gefährlich wirken} & [To not be taken care of] can be dangerous
    \end{tabular}
\end{example}
$\vm{T}(1,4)$ infinitive units can also be  used  without  \bt{zu}. Such formulations may resemble reformulated polite commands, but are completely separate from them. $\vm{T}(2,4)$ units without \bt{zu} stand in for dRAs of \bt{weil/nachdem}:
\begin{example}
    [Unordentlich geworden], sah er tot aus : Da/Nachdem er [] war, sah er tot aus\\ Ein Unfall ist beim [Nägel schneiden] passiert
\end{example}
All infinitive units are always \bt{sächlich} and \bt{einzahlig}.

\bt{*}Keep in mind that since \bt{Deutsch} does not possess the \texttt{simple} aspect, $\vm{T}_{de}(1,4)$ corresponds to  $\vm{T}_{en}(1\leftrightarrow2,4)$ in English, making both translations possible depending on context.
\subsection[Fillers]{Füllwörter}
\begin{itemize}
    \item \bt{Halt}: just. Not to be confused with the command form of \bt{halten}:\begin{example}
        Ich habe halt zu viel um meine Ohren \end{example}
    \item \bt{Eben}: just (in the context of time):\begin{example}
        Er ist eben angekommen \end{example}
    \item \bt{Ich glaube}: I think/believe
    \item \bt{Doch}:  As a conjunction `nevertheless'. As a filler, it is a stress additive without having a concrete meaning. Can be used to answer questions where the expected answer is clear from the tone and context, so \bt{doch} becomes a placeholder that answers the other person through the tone\begin{example}
        A : Wir müssen bald Geld verdienen. B : Doch (\bt{Of course})\\ Wenn ich doch mehr Geld hätte! - \bt{If [only] I had more money!}\end{example}
    \item \bt{Bitte}: Filler similar to \bt{doch}\begin{example}
        A : Meine Mutter, sie ist\ldots \\ B : Wie bitte? (\bt{What did you say?})\\ A : Ich habe gesagt, meine Mutter ist eine Maschine geworden \\ B : Bitte? (\bt{Excuse me?})\\ A : Ja, sie gehört nun den Maschinengeistern\\ B : Bitte? (\bt{What are you on?}) \end{example}
    \item \bt{Klar}: obviously/clearly/naturally
    \item \bt{Mal}: The noun meaning is `time' as in `many times, one time', but the filler meaning is varied. It mostly carries a meaning of `for the moment' or `this time' \begin{example} Probiere das mal aus \end{example}
    \item \bt{Na}: `Well, what do you think?' Behaves like \bt{doch}, but  elicits a response from the other person. It  can also be used as an informal greeting along  `Yo, what's up?' \begin{example} Na, sollen wir gehen? \end{example}
    \item \bt{Oder/ne}: Seeks a reaction to the statement at the end of which they are placed \begin{example}Du hast kein Leben, ne? \end{example}
\end{itemize}
\subsection[It]{Es}
\bt{Es} has some other uses outside of being $\vm{P}_N[a,In,3_p](1\leftrightarrow2,1)$:
\begin{enumerate}
\item Can refer to  sentence units, similar to what pronouns ($\vm{P}[a/b/c,In]$) do for noun units. References are  made to the action described in the sentence unit rather than to the entire unit:\begin{example}\begin{tabular}{m{9cm}|m{5cm}}
    Sie las ein Buch, und ich tat es auch & Ein Buch lesen\\\hline
    Die anderen waren müde. Sie war es nicht & müde sein \end{tabular}\end{example}
\item \bt{Deutsch} does not allow the case $2$ unit for an object verb to be a main sentence directly: one must  introduce such a case $2$ unit as the dRA of \bt{dass}. \bt{Es} must consequently be added as a placeholder in the ndRA as its expected object. \bt{Es} is not used if the object verb always requires main sentence case $2$ units: \begin{example}\begin{tabular}{m{14cm}}
        Ich bedauere es, dass ich nicht kommen kann\\\hline Er hat es übernommen, allen Bescheid zu geben \\\hline Ich hoffe (es), dass er reich wird \end{tabular}\end{example}
    This \bt{es} disappears if the order of dRA/ndRA is switched.
\item It is possible in general to switch the position of the units for cases $1,2$ in any main sentence (not including the dRA of a shifting relational) without consequence, assuming they are distinct entities. Since non-object verbs have no defined case $2$ unit, this switch either uses units from other cases for the purpose (if they exist), or generates an \bt{es}:\begin{example}\begin{tabular}{m{7cm}|m{7cm}}
    Ein Unfall ist an der Kreuzung geschehen & An der Kreuzung ist ein Unfall geschehen\\\hline
    Ein Unfall ist geschehen & Es ist ein Unfall geschehen \end{tabular}\end{example}
    Sometimes, this switch can use a generated \bt{es} despite the presence of other units, but the validity of such a switch is no longer guaranteed:\begin{example}
        \textcolor{red}{Es erschrickt sie vor jeder Maus} : Es ist ein Unfall an der Kreuzung geschehen \end{example}
\item If the case $1$ unit of any sentence unit is an unknown non-living thing/a natural phenomenon, \bt{es} is used to denote it:\begin{example}\begin{tabular}{c|c}
    Es klingelte plötzlich & \bt{It was ringing suddenly} \\\hline Es nieselt seit Stunden & \bt{It is drizzling since hours}\\\hline Es ist sehr spät & \bt{It is very late} \\\hline Es ist mir kalt & \bt{It is cold \ut{in me}}\end{tabular}\end{example}
    A similar use is possible when speaking about a general topic. Most of these are fixed expressions:\begin{example}\begin{tabular}{c|c}
        Es handelt sich um ($2$) & \bt{It \ut{operates itself around}} : \bt{It is about}\\\hline 
        Es hängt von ($5$) ab & \bt{It \ut{hangs from}} : \bt{It depends on}\end{tabular}\end{example}
        A preposition with its associated argument always exists in such uses. If the argument is a sentence unit, it is declared after the main sentence with a \textit{CoM} separation, and the preposition is replaced with its \bt{da}-composite (\autoref{sec:deprep}):\begin{example} Es hängt davon ab, [wie er plant, weiterzugehen]\\ Es handelt sich darum, [warum die Vögel verschwunden sind] \end{example}
\end{enumerate}
\subsubsection{Unpersönliche Leidform}
Sentences using object verbs [without an object / with non-case $2$ objects] employ [\bt{es} / the non-case $2$ object] as  the  case $1$ unit in their \bt{Leidform}:\begin{example}\begin{tabular}{c|c}
    Viele Menschen feiern & Es wird gefeiert\\\hline Sie dachten & Es war gedacht \\\hline Der Schulleiter dankte dem Förderverein & Dem Förderverein wurde gedankt\\\hline Das hilft niemandem & Nimandem ist (damit) geholfen
\end{tabular}\end{example}
In some rare circumstances, non-object verbs can form the \bt{Vorgangsleidform} in \bt{Deutsch}. I have no idea why this is allowed:
\begin{example}\begin{tabular}{c|c}
        Sie lachen & Es wird gelacht\\\hline Sie klatschen & Es wird geklatscht
\end{tabular}\end{example}
The \bt{Zustandsleidform} does not display this idiosyncrasy.

\subsection[Ancient relics]{Altdeutsch}
\begin{german}
Wenn du diese Sätze schwer verstehst, bedeutet das, du bist für diesen Schritt noch nicht bereit.
\begin{enumerate}
    \item Wesfall Eigenschaftswörter: Manchmal, anstatt den Wesfall ($6$) auf ein Nennwort anzuwenden, setzt man  ein laut $\vm{G}(4,2)$ gebeugtes Eigenschaftswort in deren Nennworteinheit ein. Solche Eigenschaftswörter  bleiben immer unverändert, unabhängig vom Stand  des verbundenen Nennworts \begin{example}
        mit voller Genuss, Erlanger Kirche, Bamberger Dom, Nürnberger Rathaus \end{example}
        Bezüglich dieser Regel sind auch einige alte Bauweisen seltens zu beobachten, Überbleibsel aus einer alten Zeit, die ihre alten Bedeutungen verloren haben, und drücken heute etwas ganz anders aus: \begin{example} \begin{tabular}{m{14cm}}Eines Tages : An diesem willkürlichen Tag, der lediglich als  Ausgangspunkt für diese Geschichte betrachtet wird\\\hline Eines Nachts : An einer beliebigen Nacht, die jetzt als Ausgangspunkt erachtet wird \end{tabular}\end{example}
    \item Redewendungen: Wenn ein Satzteil die bisher erläuterten grammatischen Regeln willkürlich durchbricht,  zusammen nicht passende Wörter besitzt, oder eine unerwartete Bedeutung aufweist,  ist er höchstwahrscheinlich  eine Redewendung \begin{example} \begin{tabular}{c}
        Kommt Zeit, kommt Rat : \bt{Things will automatically get better with time}\\\hline Im Eifer des Gefechts : \bt{In the heat of the moment} \end{tabular}\end{example}
    \item Es gibt: Eine Floskel aus einem alten Bau, in dem das Tätigkeitswort \enquote{geben} Dasein darstellte. Die Floskel ändert sich auf keinen Fall, und braucht dazu eine Fall $2$ Nennworteinheit
\end{enumerate}
\end{german}